\chapter*{}
\vspace{-3em}
{\raggedleft\huge\bfseries Conclusion générale\par}
\vspace{2em}
\label{chap:concl}

Ce rapport présente le travail accompli dans le cadre de notre projet de développement portant sur la conception et le développement de l'application Planora. Il expose en détail les différentes étapes suivies pour analyser, concevoir et réaliser une plateforme web de gestion de projet, itération par itération, en suivant une méthodologie hybride.\\
Le champ d'application de Planora englobe plusieurs modules, notamment la gestion des utilisateurs et de leurs rôles, la gestion des espaces de travail, la création et le suivi des projets, la gestion du backlog, des tâches, des sous-tâches et des sprints, ainsi que la collaboration entre les membres à travers la messagerie et les réunions. L'application intègre également un module d'intelligence artificielle permettant d'analyser les messages échangés afin de suivre la santé de l'équipe et de détecter les signes de stress ou de frustration.\\
Au cours de la réalisation de ce projet, nous avons eu l'opportunité d'appliquer et de renforcer nos connaissances en développement web, en architecture logicielle, en conception UML, en gestion de projet agile et en intégration de services intelligents. Ce travail nous a également permis de mieux comprendre les besoins liés à la collaboration au sein des équipes projet et l'importance d'une solution centralisée pour améliorer l'organisation et le suivi.\\
En ce qui concerne les perspectives futures, plusieurs améliorations peuvent être envisagées, notamment l'enrichissement du module d'intelligence artificielle, l'ajout de notifications en temps réel plus avancées, l'intégration avec des outils externes tels que GitHub, Google Calendar ou Slack, ainsi que l'amélioration des tableaux de bord analytiques. Ces évolutions permettront de renforcer l'efficacité de Planora et d'offrir une expérience plus complète aux équipes de projet.